\begingroup
\renewcommand{\Rone}{\mathcal R_1}
\section{Node CT: exact cubic tensor and retraction of CT-JM}
\label{module:CT}
\noindent\textit{Audited source: \nolinkurl{CT_complete_tensor.tex}.}
\subsection{Two functional spaces and the physical support quotient}
\label{CT:sec:functional-spaces}

The earlier working copies used the same symbol $\mathcal X$ for two
different objects.  We separate them here.  For an arbitrary boundary
vector field $W:\T\to\mathbb R^2$ let
\begin{equation}\label{CT:eq:Y-norm}
 \|W\|_{\Y}^{2}
 :=\sum_{n\in\mathbb Z}(1+|n|)|\widehat W(n)|^2.
\end{equation}
Thus $\Y=H^{1/2}(\T;\mathbb R^2)$ is the fixed space in which the disk
Taylor tensors are estimated.

The support space is a different, fan-dependent quotient.  Choose side
normals $\nu_i=e_r(\theta_i)$, put $\tau_i=e_\theta(\theta_i)$, and write
\begin{equation}\label{CT:eq:support-area-row}
 \mathsf Z_\alpha
 :=\left\{z\in\mathbb R^N:
 \sum_i\left(\tan\frac{\alpha_{i-1}}2+
              \tan\frac{\alpha_i}2\right)z_i=0\right\}
 \big/\{(q\cdot\nu_i)_i:q\in\mathbb R^2\}.
\end{equation}
For $x=\theta-\theta_i$ on
$K_i=[\theta_i-\alpha_{i-1}/2,\theta_i+\alpha_i/2]$, set
\begin{align}
 a_i^-&=\frac{z_i-z_{i-1}}{\sin\alpha_{i-1}}
          -z_i\tan\frac{\alpha_{i-1}}2,\notag\\
 a_i^+&=\frac{z_{i+1}-z_i}{\sin\alpha_i}
          +z_i\tan\frac{\alpha_i}2,\label{CT:eq:support-endpoint-jets}\\
 t_i(x)&=\frac{\tan x+\tan(\alpha_{i-1}/2)}
 {\tan(\alpha_{i-1}/2)+\tan(\alpha_i/2)}.\notag
\end{align}
The normal component of the genuine straight-side material velocity is
\begin{equation}\label{CT:eq:physical-normal-trace}
 v_\alpha(z)(\theta)
 =z_i\cos x+\bigl\{(1-t_i(x))a_i^-+t_i(x)a_i^+\bigr\}\sin x.
\end{equation}
Its tangential component, needed only in the curvature row, is
\begin{equation}\label{CT:eq:physical-tangential-trace}
 w_\alpha(z)(\theta)
 =-z_i\sin x+\bigl\{(1-t_i(x))a_i^-+t_i(x)a_i^+\bigr\}\cos x.
\end{equation}
Let $\PiNG$ delete the constant and first harmonics.  The physical Schur
form and the support Hilbert norm are, by definition,
\begin{equation}\label{CT:eq:Xalpha-definition}
 \Gtilde_\alpha(z)
 :=\qD(\PiNG v_\alpha(z)),
 \qquad
 \|z\|_{\X_\alpha}:=\Gtilde_\alpha(z)^{1/2}.
\end{equation}
The disk Bessel coercivity makes this a norm on the quotient
$\mathsf Z_\alpha$.  When a physical support slot is denoted by
$Z_\alpha(z)$, we use the unambiguous convention
\begin{equation}\label{CT:eq:support-X-convention}
 \|Z_\alpha(z)\|_{\X_\alpha}:=\|z\|_{\X_\alpha}.
\end{equation}
In particular, the comparison imported later by NR is not an additional
theorem: with the current notation it is the identity
\begin{equation}\label{CT:eq:exact-norm-interface}
 \boxed{\|Z_\alpha(z)\|_{\X_\alpha}^2
        =\Gtilde_\alpha(z).}
\end{equation}

For later bookkeeping put
$A_i(x)=(1-t_i(x))a_i^-+t_i(x)a_i^+$.  Direct rotation gives the Cartesian
identity
\begin{equation}\label{CT:eq:cartesian-support-trace}
 v_\alpha(z)e_r+w_\alpha(z)e_\theta
 =z_i\nu_i+A_i(x)\tau_i\qquad(x\in K_i).
\end{equation}
The endpoint formulas \eqref{CT:eq:support-endpoint-jets} make the two Cartesian
values at every common vertex identical.  They do \emph{not}, however, give
a ratio-free reconstruction of the full vector trace from its normal
component.

\begin{proposition}[Failure of full-vector reconstruction]
\label{CT:prop:no-vector-reconstruction}
There is no constant independent of $N$ for which
\[
 \|v_\alpha(z)e_r+w_\alpha(z)e_\theta\|_{H^{1/2}}
 \le C\|\PiNG v_\alpha(z)\|_{H^{1/2}}
\]
holds on all exact support classes.
\end{proposition}

\begin{proof}
Take an even regular fan, $\alpha_i=a=2\pi/N$, and the alternating support
mode $z_i=(-1)^i$, which satisfies the area row and has no translation
component.  The endpoint formulas give
\[
 a_i^-=z_i\cot(a/2),\qquad a_i^+=-z_i\cot(a/2).
\]
Thus the normal trace has a fixed normalized-cell profile of amplitude
$O(1)$, while the tangential coefficient $A_i$ runs between values of size
$\cot(a/2)\asymp a^{-1}$.  Scale invariance of the cell
$\dot H^{1/2}$ seminorm and summation over the $N$ cells give
\[
 \|\PiNG v_\alpha(z)\|_{H^{1/2}}^2\asymp N,
 \qquad
 \|v_\alpha(z)e_r+w_\alpha(z)e_\theta\|_{H^{1/2}}^2
 \asymp Na^{-2}.
\]
The squared ratio therefore diverges like $a^{-2}$.
\end{proof}

Consequently tangential material velocities may enter only through the
completed physical-coordinate chain rule in \eqref{CT:eq:physical-chain}; they
cannot be treated as an $H^{1/2}$ support slot.  This is exactly the
vertex-reparametrization issue isolated below.

\subsection{Fixed-domain coefficient ledger}

Let $E$ be a fixed annular extension from boundary vector fields to
$W^{1,\infty}(\mathbb D;\mathbb R^2)$, vanishing on $\{|x|\le1/2\}$.
For three boundary directions $V_i$ put
\[
 F_t(x)=x+\sum_{i=1}^3t_iEV_i(x),\qquad
 G_i=D(EV_i),\qquad G(t)=\sum_it_iG_i.
\]
Write
\begin{equation}\label{CT:eq:JC}
 J(t)=\det(I+G(t)),\qquad
 P(t)=(I+G(t))^{-1},\qquad
 C(t)=J(t)P(t)P(t)^T.
\end{equation}
If $I$ is an ordered list of distinct indices, a subscript $I$ denotes
$\partial_I|_{t=0}$, and $I_0\dot\cup\cdots\dot\cup I_r=I$ denotes an
ordered set partition, empty blocks being allowed.

All coefficient derivatives through cubic order are explicit.  Namely,
\begin{align}
 J_i&=\operatorname {tr}G_i,\notag\\
 J_{ij}&=(\operatorname {tr}G_i)(\operatorname {tr}G_j)
          -\operatorname {tr}(G_iG_j),\qquad J_{ijk}=0,
 \label{CT:eq:J-jets}\\
 P_i&=-G_i,\notag\\
 P_{ij}&=G_iG_j+G_jG_i,\notag\\
 P_{ijk}&=-\sum_{\pi\in S_3}
 G_{\pi(i)}G_{\pi(j)}G_{\pi(k)},
 \label{CT:eq:P-jets}\\
 C_I&=\sum_{I_0\dot\cup I_1\dot\cup I_2=I}
 J_{I_0}P_{I_1}P_{I_2}^{T}.
 \label{CT:eq:C-jets}
\end{align}
Here $J_\varnothing=1$ and $P_\varnothing=I$.  Formula
\eqref{CT:eq:C-jets}, rather than an ellipsis, is the metric/Jacobian ledger.

On $H_0^1(\mathbb D)$ define
\begin{align}
 A_t(u,v)&=\int_{\mathbb D}\nabla u\cdot C(t)\nabla v,
 &M_t(u,v)&=\int_{\mathbb D}J(t)uv,\label{CT:eq:forms}\\
 s(t)&=\pi^{-1}\int_{\mathbb D}J(t),
 &\widehat A_t&=s(t)A_t.\label{CT:eq:scale-form}
\end{align}
Thus the generalized eigenvalue of $(\widehat A_t,M_t)$ is the first
Dirichlet eigenvalue after normalization to area $\pi$.  Its scale rows are
\begin{equation}\label{CT:eq:scale-jets}
 \widehat A_I
 =\sum_{I_0\dot\cup I_1=I}s_{I_0}A_{I_1},
 \qquad
 A_I(u,v)=\int\nabla u\cdot C_I\nabla v,
 \qquad M_I(u,v)=\int J_Iuv.
\end{equation}
No area constraint has yet been imposed in \eqref{CT:eq:scale-jets}; in
particular, the scale terms have not been silently deleted.

To use a fixed inner product, let $S(t)$ be multiplication by
$J(t)^{-1/2}$ and set
\begin{equation}\label{CT:eq:H}
 H(t)=S(t)\widehat A_tS(t)
\end{equation}
in the form sense on $L^2(\mathbb D)$.  The mass/Jacobian jets are
\begin{align}
 S_i&=-\tfrac12J_i,\notag\\
 S_{ij}&=\tfrac34J_iJ_j-\tfrac12J_{ij},\notag\\
 S_{ijk}&=-\tfrac{15}{8}J_iJ_jJ_k
 +\tfrac34(J_{ij}J_k+J_{ik}J_j+J_{jk}J_i),
 \label{CT:eq:S-jets}
\end{align}
and the completed operator jets are
\begin{equation}\label{CT:eq:H-jets}
 H_I=\sum_{I_0\dot\cup I_1\dot\cup I_2=I}
 S_{I_0}\widehat A_{I_1}S_{I_2}.
\end{equation}
Equations \eqref{CT:eq:J-jets}--\eqref{CT:eq:H-jets} list every direct metric,
Jacobian, mass, and scale monomial through third order.

\subsection{The full symmetric state tensor}

Let $u_0$ be the positive normalized disk state, $H_0u_0=\lambda_0u_0$,
and put
\begin{equation}\label{CT:eq:R}
 Q=I-|u_0\rangle\langle u_0|,
 \qquad R=Q(H_0-\lambda_0)^{-1}Q,
 \qquad h_i=QH_iu_0.
\end{equation}
The inverse in \eqref{CT:eq:R} is positive because $\lambda_0$ is the first
simple eigenvalue.

\begin{proposition}[Complete third reduced-resolvent formula]
\label{CT:prop:full-tensor}
For arbitrary directions $V_1,V_2,V_3$, the full symmetric third
derivative is
\begin{align}
 \lambda_{123}={}&\langle u_0,H_{123}u_0\rangle\label{CT:eq:full-tensor}\\
 &-\sum_{\rm cyc(ij,k)}
 \bigl\{\langle QH_{ij}u_0,Rh_k\rangle
          +\langle h_k,RQH_{ij}u_0\rangle\bigr\}\notag\\
 &+\sum_{\pi\in S_3}
 \left\langle h_{\pi(1)},
 R(H_{\pi(2)}-\lambda_{\pi(2)}I)R
 h_{\pi(3)}\right\rangle.\notag
\end{align}
Here $\lambda_i=\langle u_0,H_iu_0\rangle$.  At the scale-normalized disk
$\lambda_i=0$ for every boundary direction.
\end{proposition}

\begin{proof}[Proof of Proposition~\ref{CT:prop:full-tensor}]
Choose the analytic normalized eigenvector $u(t)$ with
$\langle u(t),u(t)\rangle=1$.  The first differentiated state equation
gives
\begin{equation}\label{CT:eq:first-state}
 Qu_i=-Rh_i,
 \qquad \langle u_i,u_0\rangle=0.
\end{equation}
Differentiate once more, project by $Q$, and use the differentiated
normalization to eliminate the component along $u_0$.  Substitution in
$\lambda_i=\langle u,H_iu\rangle$ gives \eqref{CT:eq:full-tensor}.
For a repeated direction it reduces to
\[
 \lambda'''=\langle H'''\rangle
 -6\langle QH''u_0,Rh\rangle
 +6\langle h,R(H'-\lambda'I)Rh\rangle,
\]
which also checks all multiplicities.  Polarization proves the displayed
symmetric formula.  Criticality follows from dilation-normalized disk
stationarity.
\end{proof}

\begin{proposition}[Two energetic slots in the disk cubic tensor]
\label{CT:prop:cubic-tame}
The complete tensor of Proposition~\ref{CT:prop:full-tensor} extends to two
$\Y$ slots and one Lipschitz coefficient slot, and
\begin{equation}\label{CT:eq:cubic-tame}
 |D^3\mathcal L(0)[U,V,W]|
 \le C\sum_{\rm cyc}
 \|U\|_\Y\|V\|_\Y\|W\|_{W^{1,\infty}}.
\end{equation}
The statement concerns the complete vector tensor; no radial
reparametrization row has been separated from it.
\end{proposition}

\begin{proof}
For smooth traces use the exact physical chain rule
\eqref{CT:eq:physical-chain}.  The normal tensor is the null form
\eqref{CT:eq:null} plus the seven order-one rows
\eqref{CT:eq:seven-rows}.  Its multipliers have total order at most two and
obey the stated allocation by the dyadic argument of
Lemma~\ref{CT:lem:seven-row-audit}.

It remains to audit the curvature part.  The disk Hessian has symbol
$d^2(1+\mu_n)$ on $|n|\ge2$, hence order one with the standard discrete
differences.  The three terms in $\Gamma_{ij}$ have orders zero, one, and
one.  After symmetric polarization, their multipliers are finite
permutations of
\begin{equation*}
 (1+\mu_{n_1}),\qquad
 n_2(1+\mu_{n_1}),
 \qquad n_1+n_2+n_3=0,
\end{equation*}
up to the harmless shifts by one caused by $e_r,e_\theta$.  They have
total order at most two.  On each dyadic block put the two largest half
derivatives in $L^2$ and the remaining coefficient, together with its one
allowed derivative, in $L^\infty$.  The finite-rank area and translation
rows vanish by \eqref{CT:eq:geometric-kernel}.  Discrete Coifman--Meyer
summation proves \eqref{CT:eq:cubic-tame}.  Density defines the two energetic
slots and preserves extension independence.
\end{proof}

The uncancelled rows in \eqref{CT:eq:full-tensor} are now literally:
\begin{center}
\begin{longtable}{>{\raggedright\arraybackslash}p{.23\linewidth}
                  >{\raggedright\arraybackslash}p{.66\linewidth}}
\toprule
row & exact source\\
\midrule
\endfirsthead
\toprule
row & exact source\\
\midrule
\endhead
third direct row & $\langle H_{123}u_0,u_0\rangle$ with
\eqref{CT:eq:H-jets};\\
second-state rows & the six $H_{ij}$--$R$--$H_k$ pairings in the second
line of \eqref{CT:eq:full-tensor};\\
three first-state rows & the six $H_i$--$R$--$H_j$--$R$--$H_k$ pairings
in the last line;\\
state normalization & the three $-\lambda_iH_j$ insertions in that same
last line;\\
metric/Jacobian & the $C_I$ terms of \eqref{CT:eq:C-jets};\\
mass & the nonempty $S_I$ terms of \eqref{CT:eq:H-jets};\\
scale/area factor & the nonempty $s_I$ terms of
\eqref{CT:eq:scale-jets}.\\
\bottomrule
\end{longtable}
\end{center}
Thus the word ``state rows'' does not conceal a differentiated state, and
the geometric rows are finite set-partition sums.

\subsection{Physical coordinates, curvature, area, and translations}

Write a boundary vector direction as
\begin{equation}\label{CT:eq:components}
 V_i=f_i e_r+w_i e_\theta.
\end{equation}
The radial graph of the image of
$e_r+\sum t_iV_i$ has first jet $f_i$ and second jet
\begin{equation}\label{CT:eq:Gamma}
 \Gamma_{ij}=w_iw_j-w_if_j'-w_jf_i'.
\end{equation}
Indeed, take the modulus and argument of
$e^{i\theta}(1+\sum t_if_i+i\sum t_iw_i)$ and then express the modulus in
the new angular variable.  This gives \eqref{CT:eq:Gamma} before using any
invariance.

The exact area-and-centering chart adds a unique $\chi_{ij}$ in the span
of $1,\cos\theta,\sin\theta$ to the second radial jet.  Its coefficients
are
\begin{align}
 \chi_{ij}={}&c_{ij}^0+c_{ij}^1\cos\theta+c_{ij}^2\sin\theta,
 \label{CT:eq:chi}\\
 c_{ij}^0={}&-\frac1{2\pi}\int_\T
    (\Gamma_{ij}+f_if_j),\notag\\
 c_{ij}^1={}&-\frac1\pi\int_\T
    (\Gamma_{ij}+2f_if_j)\cos\theta,\notag\\
 c_{ij}^2={}&-\frac1\pi\int_\T
    (\Gamma_{ij}+2f_if_j)\sin\theta.\notag
\end{align}
These are obtained by twice differentiating
\[
 \tfrac12\int R^2=\pi,
 \qquad \tfrac13\int R^3(\cos\theta,\sin\theta)=0.
\]
They are the area row and the two translation-projection rows; there is no
unspecified finite-rank correction.

Let $\Lambda^{\rm n}(\rho)$ denote the scale-normalized scalar in the
normal radial chart.  Since $D\Lambda^{\rm n}(0)=0$, the exact cubic chain
rule is
\begin{multline}\label{CT:eq:physical-chain}
 D^3\mathcal L(0)[V_1,V_2,V_3]
 =D^3\Lambda^{\rm n}(0)[f_1,f_2,f_3]\\
 +\sum_{\rm cyc(ij,k)}D^2\Lambda^{\rm n}(0)
 [f_k,\Gamma_{ij}+\chi_{ij}].
\end{multline}
Scale and translation invariance imply
\begin{equation}\label{CT:eq:geometric-kernel}
 D^2\Lambda^{\rm n}(0)[g,\psi]=0,
 \qquad \psi\in\operatorname {span}\{1,\cos\theta,\sin\theta\}.
\end{equation}

The disk Hessian used throughout this node is the following literal
Fourier--Bessel form.  If
$d=\partial_ru_0(1)$ for the positive $L^2$-normalized disk state and
$\mathcal D_\circ$ is the multiplier in \eqref{CT:eq:DtN}, set
\begin{equation}\label{CT:eq:qD-definition}
 \qD(g,h):=d^2\int_\T
  (\PiNG g)(\mathcal D_\circ+1)(\PiNG h),
 \qquad \qD(g):=\qD(g,g).
\end{equation}
Thus
\begin{equation}\label{CT:eq:disk-Hessian}
 \boxed{\frac12D^2\Lambda^{\rm n}(0)[g,h]=\qD(g,h)},
 \qquad
 c\|\PiNG g\|_{H^{1/2}}^2
 \le \qD(g)\le
 C\|\PiNG g\|_{H^{1/2}}^2.
\end{equation}
Indeed, for a mean-zero diagonal direction the second boundary row is
$u_2|=d(f\mathcal D_\circ f+f^2/2)$.  Green's identity gives the raw
second Taylor coefficient
$d^2\int(f\mathcal D_\circ f+f^2/2)$.  Multiplication by the area factor
$1+\varepsilon^2(2\pi)^{-1}\int f^2$ adds
$j^2(2\pi)^{-1}\int f^2=(d^2/2)\int f^2$, because
$d^2=j^2/\pi$.  This proves the first identity by polarization.
Moreover, $\mu_{\pm1}=-1$ and
$\mu_n+1\asymp1+|n|$ for $|n|\ge2$, which proves the coercivity statement
and the geometric kernel \eqref{CT:eq:geometric-kernel}.

Consequently all $\chi_{ij}$ rows cancel exactly.  The surviving
$D^2\Lambda^{\rm n}[f_k,\Gamma_{ij}]$ rows are the curvature rows.  For
$V_1=V_2=fe_r$ and $V_3=we_\theta$, \eqref{CT:eq:physical-chain} becomes
\begin{equation}\label{CT:eq:tangent-special}
 D^3\mathcal L(0)[fe_r,fe_r,we_\theta]
 =-2D^2\Lambda^{\rm n}(0)[f,wf'],
\end{equation}
or, equivalently, $-4\qD(f,wf')$.  After multiplication by the Taylor
factor $1/2$ appropriate to the differentiated cubic term
$3\mathcal T_3$, the same row is $-2\qD(f,wf')$.  This distinction fixes
the multiplicity used downstream.

\begin{proposition}[Retraction of the curvature-as-remainder estimate]
\label{CT:prop:curvature-counterexample}
Let $f_\alpha$ be the normal fan field and let
$Z_\alpha(z)=v_\alpha(z)e_r+w_\alpha(z)e_\theta$ be a physical support
slot.  The formerly asserted estimate
\begin{equation}\label{CT:eq:curvature-absolute}
 \left|4\qD(f_\alpha,w_\alpha(z)f_\alpha')\right|
 \le CS^{5/2}\|Z_\alpha(z)\|_{\X_\alpha}.
\end{equation}
is false uniformly as $N\to\infty$.  Consequently the curvature row cannot
be placed by itself in the $o(1)$ mixed remainder.
\end{proposition}

\begin{proof}
Take $N=3M$, $a=2\pi/N$, and the three-periodic fan
$\alpha_{3j+r}=ap_r$ with $p=(1/6,8/3,1/6)$.  Choose a three-periodic
support vector satisfying the exact area row and the two translation rows;
after scaling it tends to $(17,-2,0)$.  In the material coordinate
$y=\theta/a$ one has
\[
 f_\alpha=a^2F_p+O(a^4),\qquad
 w_\alpha(z)f_\alpha'=G_{p,z}+O(a^2),\qquad
 \|Z_\alpha(z)\|_{\X_\alpha}^2\asymp a^{-1}.
\]
The high-frequency expansion
$\mathcal D_\circ+1=|D|+1+O(|D|^{-1})$ gives
\[
 \qD(f_\alpha,w_\alpha(z)f_\alpha')
 =2\pi d^2m_{p,z}a+O(a^2),
 \qquad m_{p,z}\ne0.
\]
Here $m_{p,z}$ is the mean Fourier pairing on the material period of length
three.  The factor $2\pi$ comes from its $2\pi/(3a)$ repetitions and the
change $d\theta=a\,dy$.  For a check independent of tail monotonicity, put
\[
 \widehat H(n)=\frac13\int_0^3H(y)e^{-2\pi iny/3}\,dy
\]
for every three-periodic material profile, and put
\[
 S_K=2\sum_{n=1}^K\frac{2\pi n}{3}
 \operatorname {Re}\{\overline{\widehat F_p(n)}
                         \widehat G_{p,z}(n)\}.
\]
On each of the three cells the two functions are polynomials of degree at
most two.  For $omega=2\pi n/3$ their integrals are evaluated by
\[
 I_0(l,u;\omega)=\frac{e^{-i\omega l}-e^{-i\omega u}}{i\omega},
 \qquad
 I_j(l,u;\omega)=
 \left[-\frac{x^je^{-i\omega x}}{i\omega}\right]_{l}^{u}
 +\frac{j}{i\omega}I_{j-1}(l,u;\omega).
\]
Substitution of the rational cell endpoints and coefficients, with
outward rational enclosures for $\pi$, sine, and cosine, gives at $K=4096$
\[
 2.23372<S_{4096}<2.23374.
\]
Moreover $F_p'$ has total variation $6$ and direct differentiation of the
three quadratic pieces, including their endpoint jumps, gives
$\operatorname {TV}(G_{p,z})<210$.  Twice integrating by parts in
$\widehat F_p$ and once in $\widehat G_{p,z}$ therefore yields
\[
 |m_{p,z}-S_K|
 \le\frac{\operatorname {TV}(F_p')\operatorname {TV}(G_{p,z})}
          {2\pi^2K},
 \qquad
 |m_{p,z}-S_{4096}|<0.01559.
\]
Thus $2.218<m_{p,z}<2.250$, in particular $m_{p,z}\ne0$.  The first ten
modes give the additional regression value $2.159716$, but are not used
for the rigorous tail separation.  The left side of
\eqref{CT:eq:curvature-absolute} is therefore of order $a$, whereas its right
side is $O(a^2)$.  This contradicts a uniform constant.
\end{proof}

The distributional curvature pairing remains part of the exact identity
\begin{multline}\label{CT:eq:curvature-as-difference}
 -4\qD(f_\alpha,w_\alpha(z)f_\alpha')\\
 =D^3\mathcal L(0)[B_\alpha,B_\alpha,Z_\alpha(z)]
  -D^3\Lambda^{\rm n}(0)
    [f_\alpha,f_\alpha,v_\alpha(z)].
\end{multline}
but Proposition~\ref{CT:prop:no-vector-reconstruction} and
Proposition~\ref{CT:prop:curvature-counterexample} show that its two terms may
not be estimated separately.  In fact even their sum is not a small
remainder, as proved below; it must be retained in a joint principal block.

\subsection{Extraction of the unique order-two symbol}

For real periodic functions define the same symmetric null form used by
the quadratic-bubble node,
\begin{equation}\label{CT:eq:null}
 \Cnull(f_1,f_2,f_3)
 =\frac13\sum_{\rm cyc}
 \int_\T f_1\{(\Dop f_2)(\Dop f_3)-f_2'f_3'\}.
\end{equation}
With $j=j_{0,1}$, the fixed nonzero disk solvability constant in our
integral convention is
\begin{equation}\label{CT:eq:kappa}
 \boxed{\kappa_{\mathbb D}=-\frac{6j^2}{\pi}}.
\end{equation}
Indeed, if $d=\partial_ru_0(1)$, then
$\partial_\lambda u_0(1)=d/(2j^2)$; the coefficient of the cubic boundary
solvability row is therefore $-2j^2$ times its angular mean.  The minus
sign comes from solving
\(
 (d/(2j^2))\,\delta\lambda+d\,R=0
\)
for the spectral increment.  Conversion from the Taylor coefficient to
$D^3$ supplies the factor $6$.  This sign is fixed by the outward radial
convention in \eqref{CT:eq:components}; it is not absorbed into the definition
of \(\Cnull\).

On the nongeometric subspace let
\begin{align}
 \mathcal D_\circ e^{in\theta}&=\mu_ne^{in\theta}\quad(n\ne0),
 &\mathcal D_\circ1&=0,\label{CT:eq:DtN}\\
 \mu_n&=j\frac{J_{|n|}'(j)}{J_{|n|}(j)}=|n|+r_n,
 &|\Delta^kr_n|&\le C_k(1+|n|)^{-1-k}.\notag
\end{align}
The Bessel equation also gives the exact second normal derivative
\begin{equation}\label{CT:eq:rr-exact}
 \partial_r^2\mathcal E_jg|_{r=1}
 =\Dop^2g+\mathcal B_1g,
 \qquad
 \mathcal B_1e^{in\theta}=(-\mu_n-j^2)e^{in\theta},
\end{equation}
where $\mathcal B_1$ has order one.

\begin{proposition}[Exact seven-row boundary formula]
\label{CT:prop:seven-row}
Let $f$ be real and $\PiZero f=f$, where $\PiZero$ deletes only the
constant mode.  Then
\begin{equation}\label{CT:eq:exact-boundary-cubic}
 D^3\Lambda^{\rm n}(0)[f,f,f]
 =\kappa_{\mathbb D}\,\mathfrak B(f),
\end{equation}
where
\begin{multline}\label{CT:eq:B-exact}
 \mathfrak B(f)
 =\int_\T f\,\mathcal D_\circ
       \left(f\mathcal D_\circ f+\frac12f^2\right)
 -\frac12\int_\T f^2(\Dop^2+\mathcal B_1)f\\
 +\frac{2-j^2}{6}\int_\T f^3.
\end{multline}
Consequently
\begin{equation}\label{CT:eq:seven-row}
 \mathfrak B(f)=\Cnull(f,f,f)+\sum_{a=1}^7\mathfrak r_a(f),
\end{equation}
with the seven literal lower rows
\begin{align}
 \mathfrak r_1(f)&=\int f\Dop(f\,r(D)f),
 &\mathfrak r_2(f)&=\int f\,r(D)(f\Dop f),\\
 \mathfrak r_3(f)&=\int f\,r(D)(f\,r(D)f),
 &\mathfrak r_4(f)&=\frac12\int f\Dop(f^2),\\
 \mathfrak r_5(f)&=\frac12\int f\,r(D)(f^2),
 &\mathfrak r_6(f)&=-\frac12\int f^2\mathcal B_1f,
 \label{CT:eq:seven-rows}\\
 \mathfrak r_7(f)&=\frac{2-j^2}{6}\int f^3.&&
\end{align}
Here $r(D)=\mathcal D_\circ-\Dop$ on nonzero modes and is zero on the
constant mode.  The polarized formulas are obtained by symmetric
polarization of \eqref{CT:eq:seven-rows}.
\end{proposition}

\begin{proof}
Write $u_0(r)=c_0J_0(jr)$, $d=\partial_ru_0(1)$, and expand without
factorials
\begin{equation*}
 u=u_0+\varepsilon u_1+\varepsilon^2u_2+
   \varepsilon^3u_3+O(\varepsilon^4),
 \qquad
 \lambda=j^2+\varepsilon^2\lambda_2+
   \varepsilon^3\lambda_3+O(\varepsilon^4).
\end{equation*}
The linear term vanishes because $f$ has zero mean.  The first two
Dirichlet rows are exactly
\begin{align}
 u_1|_{r=1}&=-df,
 &u_2|_{r=1}
 &=d\left(f\mathcal D_\circ f+\frac12f^2\right).
 \label{CT:eq:first-two-boundary-rows}
\end{align}
The nonzero angular modes of $\partial_ru_2$ are therefore
$d\mathcal D_\circ(f\mathcal D_\circ f+f^2/2)$.  Its constant mode contains
the second spectral correction and the state-normalization correction, but
it disappears after pairing with the mean-zero function $f$.  The Bessel
equation gives the exact identities
\begin{equation}\label{CT:eq:radial-identities}
 \partial_r^2u_1|_{r=1}
 =-d(\Dop^2+\mathcal B_1)f,
 \qquad
 \partial_r^3u_0(1)=(2-j^2)d.
\end{equation}

The third Dirichlet row is
\begin{equation*}
 u_3|_{r=1}
 =-f\partial_ru_2-\frac12f^2\partial_r^2u_1
   -\frac16f^3\partial_r^3u_0.
\end{equation*}
Choose the state normalization with $\langle u_1,u_0\rangle=0$.
Green's identity then gives
$\lambda_3=d\int_\T u_3|_{r=1}$.  Since the Rellich identity yields
$d^2=j^2/\pi$, multiplication by $3!=6$ gives
$-6d^2=\kappa_{\mathbb D}$.  Substitution of
\eqref{CT:eq:first-two-boundary-rows}--\eqref{CT:eq:radial-identities} proves
\eqref{CT:eq:exact-boundary-cubic}--\eqref{CT:eq:B-exact}.

Finally insert $\mathcal D_\circ=\Dop+r(D)$.  The two terms containing
only $\Dop$ are
\begin{equation*}
 \int f\Dop(f\Dop f)-\frac12\int f^2\Dop^2f
 =\Cnull(f,f,f).
\end{equation*}
The remaining terms are exactly \eqref{CT:eq:seven-rows}.
\end{proof}

\begin{lemma}[Seven-row multiplier audit]
\label{CT:lem:seven-row-audit}
For the symmetric polarization $\mathfrak r_a(f_1,f_2,f_3)$ of every row
in \eqref{CT:eq:seven-rows},
\begin{equation}\label{CT:eq:seven-row-estimate}
 |\mathfrak r_a(f_1,f_2,f_3)|
 \le C\sum_{p<q}\|f_p\|_{H^{1/2}}
                    \|f_q\|_{H^{1/2}}
                    \|f_\ell\|_{L^\infty}.
\end{equation}
The constant is uniform over $a=1,\ldots,7$.
\end{lemma}

\begin{proof}
The symbols of $r(D)$ and $\mathcal B_1$ satisfy
\begin{align*}
 |\Delta^kr_n|&\le C_k(1+|n|)^{-1-k},\\
 |\Delta^k(-\mu_n-j^2)|&\le C_k(1+|n|)^{1-k}.
\end{align*}
After polarization, the seven multipliers are finite permutations of
\begin{multline}\label{CT:eq:seven-symbols}
 |n_2|r_{n_3},\quad r_{n_2}|n_3|,\quad
 r_{n_2}r_{n_3},\quad |n_2|,\quad r_{n_2},\quad
 -\mu_{n_3}-j^2,\quad1,\\
 n_1+n_2+n_3=0.
\end{multline}
Each multiplier has total order at most one and satisfies
\begin{equation*}
 |\Delta^\gamma m(n_1,n_2,n_3)|
 \le C_\gamma(1+\max_i|n_i|)^{1-|\gamma|}.
\end{equation*}
On a dyadic block the two largest frequencies are comparable.  Put the
corresponding half derivatives in $L^2$ and the remaining low-frequency
factor in $L^\infty$.  Discrete Coifman--Meyer summation, followed by
Cauchy--Schwarz in the common high index, proves
\eqref{CT:eq:seven-row-estimate}.  The finitely many modes with
$|n|\le2$ are absorbed by the inhomogeneous norms.
\end{proof}

\begin{theorem}[Exact cubic split and literal remainder formula]
\label{CT:thm:complete-split}
For mean-zero normal directions, $\PiZero f_i=f_i$, one has the exact
identity
\begin{equation}\label{CT:eq:complete-split}
 \boxed{
 D^3\Lambda^{\rm n}(0)[f_1,f_2,f_3]
 =\kappa_{\mathbb D}\Cnull(f_1,f_2,f_3)
  +\Rone(f_1,f_2,f_3).}
\end{equation}
The remainder is not defined by an order symbol.  Its literal formula is
\begin{align}
 \Rone(f_1,f_2,f_3)={}&
 \langle u_0,H_{123}u_0\rangle
 -\sum_{\rm cyc(ij,k)}
 \bigl\{\langle QH_{ij}u_0,Rh_k\rangle
        +\langle h_k,RQH_{ij}u_0\rangle\bigr\}
 \notag\\
 &+\sum_{\pi\in S_3}
 \left\langle h_{\pi(1)},
 R(H_{\pi(2)}-\lambda_{\pi(2)}I)R
 h_{\pi(3)}\right\rangle
 -\kappa_{\mathbb D}\Cnull(f_1,f_2,f_3),
 \label{CT:eq:R-literal}
\end{align}
where every $H_I$ is the finite set-partition expression
\eqref{CT:eq:H-jets}, with \eqref{CT:eq:scale-jets}, \eqref{CT:eq:C-jets}, and
\eqref{CT:eq:J-jets} substituted.  Thus \eqref{CT:eq:R-literal}, unlike the old
tag names alone, specifies every scalar monomial and its multiplicity.
Equivalently, after the necessary direct/state cancellations have been
performed, the same remainder is the seven-row grouped expression
\begin{equation}\label{CT:eq:R-seven-polarized}
 \boxed{\Rone(f_1,f_2,f_3)
 =\kappa_{\mathbb D}\sum_{a=1}^7
   \mathfrak r_a(f_1,f_2,f_3),}
\end{equation}
where $\mathfrak r_a(f_1,f_2,f_3)$ denotes symmetric polarization of the
corresponding line of \eqref{CT:eq:seven-rows}.  Formula
\eqref{CT:eq:R-seven-polarized}, rather than the uncancelled raw monomials, is
the row decomposition used for analytic estimates.

For bookkeeping, expand \eqref{CT:eq:R-literal} fully and assign each scalar
monomial its first applicable tag in the ordered list
\begin{equation}\label{CT:eq:R-tag-priority}
 r(D),\quad \mathcal B_1,\quad C_I\ (I\ne\varnothing),\quad
 S_I\ (I\ne\varnothing),\quad s_I\ (I\ne\varnothing).
\end{equation}
This gives the following disjoint ledger:
\begin{equation}\label{CT:eq:R-ledger}
 \Rone=R_{\rm state}^{r}+R_{rr}+R_{\rm met}
       +R_{J/M}+R_{\rm scale}.
\end{equation}
Here:
\begin{enumerate}
\item $R_{\rm state}^{r}$ consists of the state monomials in
\eqref{CT:eq:full-tensor} containing at least one $r_n$ from
\eqref{CT:eq:DtN};
\item $R_{rr}$ consists of those containing $\mathcal B_1$ from
\eqref{CT:eq:rr-exact};
\item $R_{\rm met}$ consists of the remaining nonprincipal monomials
tagged by a nonempty $C_I$ in \eqref{CT:eq:C-jets};
\item $R_{J/M}$ consists of the remaining monomials whose first tag is a
nonempty $S_I$ in \eqref{CT:eq:H-jets};
\item $R_{\rm scale}$ consists of the remaining monomials whose first tag
is a nonempty $s_I$ in \eqref{CT:eq:scale-jets};
\end{enumerate}
The priority \eqref{CT:eq:R-tag-priority} makes the list disjoint, and the
set-partition formulae make it exhaustive.
\end{theorem}

\begin{proof}[Proof of Theorem~\ref{CT:thm:complete-split}]
Proposition~\ref{CT:prop:seven-row} proves the diagonal identity
\eqref{CT:eq:complete-split} together with
\eqref{CT:eq:R-seven-polarized}; symmetric polarization gives arbitrary
nongeometric slots.  Independently, Proposition~\ref{CT:prop:full-tensor}
gives the complete reduced-resolvent expression.  Subtracting the same
principal null form from that expression gives \eqref{CT:eq:R-literal}.
Hence the fixed-domain formula and the seven grouped boundary rows are two
exact representations of the same trilinear remainder.

Expanding \eqref{CT:eq:R-literal} and applying
\eqref{CT:eq:R-tag-priority} records the origin of every raw monomial in
\eqref{CT:eq:R-ledger}.  No estimate is taken before the raw metric and state
rows have been recombined into the seven rows
\eqref{CT:eq:R-seven-polarized}; this is the required order-two
direct/resolvent cancellation.
\end{proof}

\begin{corollary}[Order-one normal remainder]
\label{CT:cor:R-estimate}
For mean-zero normal slots,
\begin{equation}\label{CT:eq:R-estimate}
 |\Rone(f_1,f_2,f_3)|
 \le C\sum_{p<q}\|f_p\|_{H^{1/2}}
                    \|f_q\|_{H^{1/2}}
                    \|f_\ell\|_{L^\infty},
 \qquad \{\ell\}=\{1,2,3\}\setminus\{p,q\}.
\end{equation}
\end{corollary}

\begin{proof}
Use \eqref{CT:eq:R-seven-polarized} and sum
\eqref{CT:eq:seven-row-estimate} over the seven rows.
\end{proof}

For later fan inputs the constant scale mode must be retained explicitly.
For arbitrary normal slots write
\begin{equation*}
 g_i=\PiZero f_i,
 \qquad \psi_i=(I-\PiZero)f_i,
\end{equation*}
and, for $\varnothing\ne A\subset\{1,2,3\}$, let
$\xi_i^A=\psi_i$ if $i\in A$ and $\xi_i^A=g_i$ otherwise.  Define the
finite seven-term geometric remainder by
\begin{equation}\label{CT:eq:Rgeo}
 \Rgeo(f_1,f_2,f_3)
 :=\sum_{\varnothing\ne A\subset\{1,2,3\}}
 D^3\Lambda^{\rm n}(0)
 [\xi_1^A,\xi_2^A,\xi_3^A]
\end{equation}
and set
\begin{equation}\label{CT:eq:Rfull}
 \Rfull(f_1,f_2,f_3)
 :=\Rone(g_1,g_2,g_3)+\Rgeo(f_1,f_2,f_3).
\end{equation}
Multilinearity gives, with no omitted low mode,
\begin{equation}\label{CT:eq:normal-full-split}
 \boxed{
 D^3\Lambda^{\rm n}(0)[f_1,f_2,f_3]
 =\kappa_{\mathbb D}\Cnull(g_1,g_2,g_3)
  +\Rfull(f_1,f_2,f_3).}
\end{equation}
Every term of \eqref{CT:eq:Rgeo} contains a constant.
Apply Proposition~\ref{CT:prop:cubic-tame} with that smooth finite-mode slot
as the Lipschitz coefficient slot and the other two as energetic slots.
Together with Corollary~\ref{CT:cor:R-estimate}, this proves the same
order-one bound for $\Rfull$ as in \eqref{CT:eq:R-estimate}, with every $f_i$
replaced by its full inhomogeneous norm.  For the fan input,
\begin{equation}\label{CT:eq:fan-geometric-size}
 \|(I-\PiZero)f_\alpha\|_{W^{1,\infty}}\le CS^2,
\end{equation}
so the geometric pure and mixed rows have the lower majorants $CS^5$ and
$CS^{7/2}\|Z\|_\X$, respectively.

The physical curvature is a separate tensor, not a summand of the normal
remainder.  Define
\begin{equation}\label{CT:eq:K-physical}
 \mathcal K_{\rm curv}(V_1,V_2,V_3)
 :=\sum_{\rm cyc(ij,k)}D^2\Lambda^{\rm n}(0)
     [f_k,\Gamma_{ij}].
\end{equation}
Set
\begin{equation}\label{CT:eq:Rphys}
 \Rphys(V_1,V_2,V_3)
 :=\Rfull(f_1,f_2,f_3)
   +\mathcal K_{\rm curv}(V_1,V_2,V_3).
\end{equation}
Then \eqref{CT:eq:physical-chain} and the geometric kernel give the typed
identity
\begin{equation}\label{CT:eq:physical-complete-split}
 \boxed{
 D^3\mathcal L(0)[V_1,V_2,V_3]
 =\kappa_{\mathbb D}\Cnull(f_1,f_2,f_3)
  +\Rphys(V_1,V_2,V_3).}
\end{equation}
The $\chi_{ij}$ rows are absent from this formula because they vanish
exactly, not because they have been assigned to $\Rone$.

For clarity, the cancellations used above are recorded separately:
\begin{center}
\begin{tabular}{>{\raggedright\arraybackslash}p{.27\linewidth}
                >{\raggedright\arraybackslash}p{.61\linewidth}}
\toprule
family & cancellation or surviving order\\
\midrule
two principal state derivatives & self-adjointness and one integration by
parts combine them into \eqref{CT:eq:null};\\
mixed-sign order-two interactions & cancel inside
$(\Dop f_2)(\Dop f_3)-f_2'f_3'$;\\
direct metric versus second state & their order-two pieces reconstruct the
boundary acceleration in Proposition~\ref{CT:prop:seven-row}; after the
combination only the seven order-one rows \eqref{CT:eq:seven-rows} remain;\\
state normalization & zero because $\lambda_i=0$ at the normalized disk;\\
mass and scale & no tangential multiplier, hence order zero;\\
curvature & total order two in general; its ambient two-energetic-slot bound
does not imply a radial support bound.  Both the separate remainder estimate
and the CT--JM small-sum estimate are false; the row is retained inside the
honest JP--JC joint scalar audit;\\
area and translations & exact Hessian kernel
\eqref{CT:eq:geometric-kernel}.\\
\bottomrule
\end{tabular}
\end{center}

We next define the scalar and support functionals used in the relative
argument.  Put $\alpha_*=a\mathbf1$,
$\gamma(t)=\alpha_*+t(\alpha-\alpha_*)$, and write
$f_t=f_{\gamma(t)}$.  Fix at $t=0$ a complement to the area and two
translation rows.  Let
\begin{equation}\label{CT:eq:constraint-transport}
 C_t:\mathsf Z_{\alpha_*}\longrightarrow\mathsf Z_{\gamma(t)}
\end{equation}
be the analytic constraint transport obtained by adding the unique linear
combination of the scale and two translation vectors that satisfies the
three rows at $\gamma(t)$.  Its $3\times3$ Gram matrix is a uniform
perturbation of the regular-fan matrix on the material polydisc, so $C_t$
and $C_t^{-1}$ are uniformly bounded.  For a class $z$ in the fixed
complement define
\begin{align}
 F_1(t)&:=\Rfull(f_t,f_t,f_t),\label{CT:eq:F1-definition}\\
 \Psi_{1,n}(t)[z]
 &:=\Rfull(f_t,f_t,v_{\gamma(t)}(C_tz)),
 \label{CT:eq:Psi1n-definition}\\
\Psi_{1,t}(t)[z]
 &:=-4\qD
 (f_t,w_{\gamma(t)}(C_tz)f_t').
 \label{CT:eq:Psi1t-definition}
\end{align}
The last line is exactly the two curvature placements in
\eqref{CT:eq:physical-chain}; it is not included a second time in
$\Psi_{1,n}$.

\begin{lemma}[Material Cauchy bounds for the normal order-one rows]
\label{CT:lem:R-Cauchy}
In the uniformly comparable material regime
\eqref{CT:eq:uniformly-comparable}, $F_1$ and $\Psi_{1,n}$ extend
holomorphically to the line polydisc of radius $c/\delta(t)$, where
\begin{equation*}
 \delta(t)=\max_i\frac{|\alpha_i-a|}{\gamma_i(t)}.
\end{equation*}
Their scalar and dual majorants are
\begin{equation}\label{CT:eq:R-complex-majorants}
 |F_1|\le CS^5,\qquad
 \|\Psi_{1,n}\|_{\X_{\gamma(t)}^*}\le CS^{7/2}.
\end{equation}
Consequently, for
$E(t)=\sum_i\gamma_i(t)^2(\alpha_i-a)^2$,
\begin{align}
 |F_1''(t)|&\le CSE(t),\label{CT:eq:F1-Cauchy}\\
 \|\Psi_{1,n}'(t)\|_{\X_{\gamma(t)}^*}
 &\le CS^{3/2}E(t)^{1/2}.\label{CT:eq:Psi1n-Cauchy}
\end{align}
\end{lemma}

\begin{proof}
Split every polar cell at its normal and pull both halves to $[0,1]$.
The secant profile, the traces
\eqref{CT:eq:physical-normal-trace}--\eqref{CT:eq:physical-tangential-trace},
and the constraint matrix in \eqref{CT:eq:constraint-transport} are rational
analytic functions of the incident widths.  Their denominators stay in a
fixed sector on the stated polydisc.  The normal estimate
\eqref{CT:eq:R-estimate} together with the finite geometric expansion
\eqref{CT:eq:Rgeo} gives both bounds in
\eqref{CT:eq:R-complex-majorants}.  Same and adjacent cells use
the normalized jets.  On separated cells, subtracting both cell means
makes the first-separating dyadic annuli summable.  No bound contains the
number of cells.

One-variable Cauchy on
$\gamma(t)+\zeta(\alpha-\alpha_*)$ gives respectively
$CS^5\delta(t)^2$ and $CS^{7/2}\delta(t)$.  If $k$ realizes $\delta(t)$, comparability gives
\begin{equation*}
 E(t)\ge\gamma_k(t)^2(\alpha_k-a)^2
       \asymp S^4\delta(t)^2.
\end{equation*}
This proves \eqref{CT:eq:F1-Cauchy}--\eqref{CT:eq:Psi1n-Cauchy}.
\end{proof}

\begin{lemma}[Relative closure of the pure cubic row]
\label{CT:lem:R-relative}
Use the material notation of Section~\ref{CT:sec:material}.  Let $f_\alpha$
be the exact normal fan input and let
$Z_\alpha(z)$ be one physical support slot.  On the exact scale/translation
quotient there is a modulus $\varepsilon_{1}\to0$ such that
\begin{equation}
 |\Rfull(f_\alpha,f_\alpha,f_\alpha)
  -\Rfull(f_{a\mathbf1},f_{a\mathbf1},f_{a\mathbf1})|
 \le\varepsilon_1(S)\Jfour(\alpha).\label{CT:eq:R-relative-pure}
\end{equation}
Equivalently, in the typed physical notation \eqref{CT:eq:Rphys}, the left
side is
\begin{equation}
 \left|\Rphys(B_\alpha,B_\alpha,B_\alpha)
       -\Rphys(B_{a\mathbf1},B_{a\mathbf1},B_{a\mathbf1})\right|.
\end{equation}
\end{lemma}

\begin{proof}
The full remainder estimate following \eqref{CT:eq:Rfull},
\eqref{CT:eq:fan-geometric-size}, and the fan sizes give absolute majorants
$CS^5$ for the pure row.

Fix $0<\beta<1$.  If $\Jfour\ge S^{5-\beta}$, these bounds divided by the
right side of \eqref{CT:eq:R-relative-pure} gives $CS^\beta$.
In the complementary regime the exact quartic Bregman identity
\eqref{CT:eq:J4-bregman-comparability} and the quantitative estimate following
it give
$\|\alpha-a\mathbf1\|_\infty/S\le2S^{(1-\beta)/2}$, so for small $S$ the
entire segment is uniformly comparable and lies in the material polydisc
\eqref{CT:eq:material-polydisc}.
Lemma~\ref{CT:lem:R-Cauchy} gives the required material derivative bound.
Cyclic symmetry kills the first pure derivative at the regular fan.
Integrate with \eqref{CT:eq:action}.  This proves
\eqref{CT:eq:R-relative-pure} with no cell-count factor.
\end{proof}

\begin{proposition}[Retraction of CT--JM]
\label{CT:prop:joint-mixed-counterexample}
There is no modulus $\varepsilon_{\rm JM}(S)\to0$ for which
\begin{equation}\label{CT:eq:false-R-relative-mixed}
 \left|\Rphys(B_\alpha,B_\alpha,Z_\alpha(z))\right|
 \le\varepsilon_{\rm JM}(S)\Jfour(\alpha)^{1/2}
       \|Z_\alpha(z)\|_{\X_\alpha}
\end{equation}
holds uniformly on the exact scale/translation quotient.
\end{proposition}

\begin{proof}
Use the three-periodic family from
Proposition~\ref{CT:prop:curvature-counterexample}.  Its seven normal rows
satisfy, by the displayed order-one multiplier estimate,
\[
 \Rfull(f_\alpha,f_\alpha,v_\alpha(z))=O(a^3),
\]
whereas
\[
 -4\qD(f_\alpha,w_\alpha(z)f_\alpha')
 =-8\pi d^2m_{p,z}a+O(a^2),\qquad m_{p,z}\ne0.
\]
For this family $S\asymp a$,
$\Jfour(\alpha)\asymp a^3$, and
$\|Z_\alpha(z)\|_{\X_\alpha}\asymp a^{-1/2}$; hence the product on the
right of \eqref{CT:eq:false-R-relative-mixed} is comparable to $a$.  Dividing
by that product gives a positive limiting lower bound, contradicting
$\varepsilon_{\rm JM}(S)\to0$.
\end{proof}

The surviving leading support forcing in the differentiated cubic Taylor
row is therefore
\begin{equation}\label{CT:eq:effective-leading-forcing}
 2\qD(f_\alpha,v_\alpha(z))
 -2\qD(f_\alpha,w_\alpha(z)f_\alpha')
 =2\qD\bigl(f_\alpha,
     v_\alpha(z)-w_\alpha(z)f_\alpha'\bigr).
\end{equation}
For a polar graph $R_\alpha$, the corresponding genuine polygon-normal
trace is
\begin{equation}\label{CT:eq:effective-normal-trace}
 v_{\alpha,{\rm eff}}(z)
 =\frac{v_\alpha(z)-(R_\alpha'/R_\alpha)w_\alpha(z)}
        {\sqrt{1+(R_\alpha'/R_\alpha)^2}}.
\end{equation}
Equations \eqref{CT:eq:effective-leading-forcing}--
\eqref{CT:eq:effective-normal-trace}, together with all remaining cubic rows,
are the exact input exported to the JP/JC joint scalar audit.  No estimate
in this node replaces that joint principal analysis.

\begin{remark}[Sign and geometric tests]
The Fourier symbol in \eqref{CT:eq:null} vanishes when the two differentiated
frequencies have opposite signs and equals twice their product when they
have the same sign.  A constant in any slot gives zero by Parseval.  The
full tensor, including \eqref{CT:eq:physical-chain}, vanishes on dilation and
translation paths; the principal and lower rows need not vanish there
separately.
\end{remark}

\subsection{Uniform analytic tail on material cells}
\label{CT:sec:material}

Let
\begin{equation}\label{CT:eq:fan}
 \alpha_i>0,\qquad \sum_i\alpha_i=2\pi,\qquad
 a=2\pi/N,\qquad S=a+\max_i\alpha_i,
\end{equation}
and let $B_\alpha$ be the physical fan boundary field.

\begin{theorem}[Unconditional fixed-chart and natural-order facts]
\label{CT:thm:fixed-chart-unconditional}
On a fixed sufficiently small real or complex $W^{1,\infty}$ ball, the
pulled-back stiffness and mass forms, their compact solution operator, and
the scale-normalized simple first-eigenvalue branch $\mathcal L(W)$ are
analytic in operator norm.  On that complexified chart its bilinear
Hessian, and hence at every real near-disk active polygon the genuine
polygon-preserving Hessian on continuous physical-side-affine fields,
satisfies
\begin{equation}\label{CT:eq:polygon-natural-Hessian}
 |D^2\mathcal L(W)[U,V]|
 \le C\|U\|_{H^{1/2}(\partial P)}
       \|V\|_{H^{1/2}(\partial P)}.
\end{equation}
Here the trace norms on the complexified chart are transported to the
underlying real boundary; these norms are uniformly equivalent throughout
the fixed chart.
The constant is independent of the number, the minimum size, and the
ratios of the sides.
\end{theorem}

\begin{proof}
For the first assertion use the annular extension in the fixed-domain
ledger.  The coefficients $J$, $P$, $C$, $S$, and $H$ in
\eqref{CT:eq:JC}--\eqref{CT:eq:H} are convergent rational power series in
$D(EW)$ on a fixed $L^\infty$ ball.  Uniform ellipticity, the disk spectral
gap, the resolvent identity, and the Riesz projection therefore give the
operator-norm analytic simple branch.  This applies directly to polygonal
boundary traces and does not round their corners.

For the second assertion, extend a physical affine-side trace by a uniform
right inverse from $H^{1/2}(\partial P;\mathbb R^2)$ to
$H^1(P;\mathbb R^2)$.  The radial bi-Lipschitz maps between the disk and
the polygons give a uniform norm for this inverse.  Convexity and the
common inner and outer radii give uniform bounds for the eigenvalue, its
spectral gap, the ground state, and its gradient.  At a supporting line
the barrier $s(D-s)$ gives
$u(x)\le C\operatorname {dist}(x,\partial P)$; a scaled interior estimate
then gives $\|\nabla u\|_\infty\le C$.

In the exact second-variation formula, the first coefficient jet is linear
and the second is quadratic in the two extension gradients.  The direct
rows are bounded by the product of their $H^1$ norms; the two state
forcings have the same $H^{-1}$ bounds, and the reduced resolvent is
uniformly coercive.  On the smaller complex ball the real part of the
pulled-back form remains uniformly coercive and the same spectral contour
separates the analytic simple branch; the identical energy estimates
therefore hold for the complex bilinear Hessian.  Polarization proves
\eqref{CT:eq:polygon-natural-Hessian}.  Approximation by Lipschitz extensions
and invariance under changes with zero boundary trace justify the use of
the energy extensions.
\end{proof}

Put
\begin{equation}\label{CT:eq:R4-definition}
 \mathcal R_4(W)=\mathcal L(W)-\mathcal L(0)
 -\frac12D^2\mathcal L(0)[W,W]
 -\frac16D^3\mathcal L(0)[W,W,W].
\end{equation}

\begin{corollary}[Unconditional tame scalar fourth remainder]
\label{CT:cor:tame-scalar-fourth}
On a fixed smaller real or complex chart ball,
\begin{equation}\label{CT:eq:tame-scalar-fourth}
 \boxed{
 |\mathcal R_4(W)|
 \le C\|W\|_{H^{1/2}}^2\|W\|_{W^{1,\infty}}^2.}
\end{equation}
The constant is independent of the number, minimum size, and ratios of the
physical sides.  This is a scalar statement and makes no assertion about
the retired support-covector row of CT--M4.
\end{corollary}

\begin{proof}
The scale-normalized disk branch has $D\mathcal L(0)=0$.  For fixed $W$ put
\[
 g(\zeta)=D^2\mathcal L(\zeta W)[W,W].
\]
Theorem~\ref{CT:thm:fixed-chart-unconditional} makes $g$ analytic for
$|\zeta|<r_0/\|W\|_{W^{1,\infty}}$ and bounds it there by
$C\|W\|_{H^{1/2}}^2$.  Cauchy's coefficient estimate, after decreasing the
chart so that this radius is greater than $2$, gives for $0\le t\le1$
\[
 |g(t)-g(0)-tg'(0)|
 \le C\|W\|_{H^{1/2}}^2\|W\|_{W^{1,\infty}}^2t^2.
\]
Twice integrating the scalar Taylor formula yields
\[
 \mathcal R_4(W)
 =\int_0^1(1-t)\{g(t)-g(0)-tg'(0)\}\,dt,
\]
and proves \eqref{CT:eq:tame-scalar-fourth}.
\end{proof}

Call a material line $h_i(t)=a+td_i$ uniformly comparable when
\begin{equation}\label{CT:eq:uniformly-comparable}
 \frac a2\le h_i(t)\le\frac{3a}{2}
 \qquad(1\le i\le N,\ 0\le t\le1).
\end{equation}
For such a line put
\[
 E(t)=\sum_i h_i(t)^2d_i^2.
\]

\begin{remark}[Retired CT--M4 proposal: not an active obligation]
\label{CT:obl:material-fourth}
Historically, this proposal assumed that CT--JP supplied a no-ratio Hilbert
form $\Gjp_\alpha$ on the exact support quotient and wrote
$\|z\|_{\Xjp_\alpha}=\Gjp_\alpha(z)^{1/2}$.  After splitting every physical
cell at its normal and pulling the two halves to fixed unit intervals,
prove the following four estimates with a constant independent of the
number, minimum size, and ratios of the cells and moduli
$\omega_4,\omega_{4,s}\downarrow0$ at the origin:
\begin{align}
 |\mathcal R_4(B_\alpha)|&\le CS^5,\notag\\
 \|D\mathcal R_4(B_\alpha)\|_{(\Xjp_\alpha)^*}&\le CS^{7/2},
 \label{CT:eq:M4-size}\\
 \left|\frac{d^2}{dt^2}\mathcal R_4(B_{h(t)})\right|
 &\le \omega_4(S) E(t),\notag\\
 \left\|\frac d{dt}\left\{D\mathcal R_4(B_{h(t)})
     \circ\mathcal T_t\right\}\right\|_{(\Xjp_\alpha)^*}
 &\le \omega_{4,s}(S)E(t)^{1/2},
 \label{CT:eq:M4-path}
\end{align}
Here the varying support fibers are identified by the uniformly bounded
constraint-preserving material transport
\[
 \mathcal T_t:\Xjp_\alpha\longrightarrow\Xjp_{h(t)},
 \qquad
 \|\mathcal T_t\|+\|\mathcal T_t^{-1}\|\le C,
 \qquad \mathcal T_1=I.
\]
It is obtained by the fixed half-cell pullback followed by the exact
area/translation projection.  In the second estimate of
\eqref{CT:eq:M4-path}, the derivative is the total derivative of the displayed
fixed-fiber covector and therefore includes the variation of the physical
support field as the fan changes.  The estimates concern the completed direct-plus-state
functional; its order-two rows may not be estimated separately.
\end{remark}

\begin{remark}[Why the old CT--M4 proposal was insufficient]
Theorem~\ref{CT:thm:fixed-chart-unconditional} proves scalar analyticity and
an absolute $H^{1/2}$ Hessian bound, but neither gives the vanishing moduli
in \eqref{CT:eq:M4-path}.  Those moduli express the cancellation between the
direct metric row and the reduced-resolvent row at the natural boundary
order.  Historical rounded Calder\'on arguments estimated those rows
separately and therefore do not prove this obligation.

Corollary~\ref{CT:cor:tame-scalar-fourth} now closes the first, purely scalar
bound in \eqref{CT:eq:M4-size}.  It does not close the support-covector bound
there, either relative path estimate in \eqref{CT:eq:M4-path}, or the moving
support-fiber comparison.  Thus it is sufficient for JP--BU but does not
resurrect CT--M4.

The former fifth ``support Hessian'' row is not part of CT--M4.  It used the
wrong radial Schur metric and cannot control the physical support problem.
Its correct replacement is the exact support Hessian in JP--EH, integrated
without a post-cubic support remainder.

The former CT--AH hypothesis---bounded holomorphy of
$W\mapsto D^2\mathcal L(W)$ in
$\operatorname {Bil}(H^{1/2}\times H^{1/2})$ on a uniform complex
$W^{1,\infty}$ ball---does not by itself imply CT--M4, because the ambient
full-vector norm need not be uniformly comparable to the joint support
metric.  CT--AH is no longer a vertex of the active DAG.
\end{remark}

The no-ratio physical sizes are
\begin{equation}\label{CT:eq:fan-sizes}
 \|B_\alpha\|_\Y^2\le CS^3,\qquad
 \|B_\alpha\|_{W^{1,\infty}}\le CS.
\end{equation}
Split each material cell at its normal and pull both halves to $[0,1]$.
The secant, support, Jacobian, and finite constraint coefficients are
rational analytic functions of the incident widths.  Hence, on a fan
satisfying \eqref{CT:eq:uniformly-comparable}, they extend to
\begin{equation}\label{CT:eq:material-polydisc}
 |\zeta_i-\alpha_i|<c\alpha_i,\qquad
 \sum_i\zeta_i=2\pi,
\end{equation}
with the same bounds as \eqref{CT:eq:fan-sizes}.  The denominator sector and
the chart radius are independent of $N$.

\begin{lemma}[Archived consequences of the retired CT--M4 proposal]
\label{CT:lem:Cauchy}
Assume the retired CT--M4 proposal \ref{CT:obl:material-fourth}.  Let
\begin{align*}
 F_4(\alpha)&=\mathcal R_4(B_\alpha),\\
 \Psi_4(\alpha)[Z]&=D\mathcal R_4(B_\alpha)[Z].
\end{align*}
At the real fan $\alpha$,
\begin{equation}\label{CT:eq:tail-majorants}
 |F_4|\le CS^5,\qquad
 \|\Psi_4\|_{(\Xjp_\alpha)^*}\le CS^{7/2}.
\end{equation}
If $h_i(t)=a+t d_i$ is uniformly comparable in the sense of
\eqref{CT:eq:uniformly-comparable} and
$E(t)=\sum_i h_i(t)^2d_i^2$, then
\begin{equation}\label{CT:eq:Cauchy-bounds}
 |\partial_t^2F_4(h(t))|\le \omega_4(S)E(t),
 \qquad
 \left\|\partial_t\{\Psi_4(h(t))\circ\mathcal T_t\}
     \right\|_{(\Xjp_\alpha)^*}
 \le \omega_{4,s}(S)E(t)^{1/2}.
\end{equation}
\end{lemma}

\begin{proof}
The two bounds in \eqref{CT:eq:tail-majorants} are
\eqref{CT:eq:M4-size}; the two bounds in \eqref{CT:eq:Cauchy-bounds} are
\eqref{CT:eq:M4-path}.  The identification of the varying dual spaces is the
fixed material pullback required in CT--M4, so no derivative of a moving
corner or hidden sum over material coordinates is introduced.
\end{proof}

The quartic action identities
\begin{equation}\label{CT:eq:action}
 \Jfour(\alpha)=12\int_0^1(1-t)E(t)\,dt,
 \qquad
 \int_0^1E(t)\,dt\le C\Jfour(\alpha)
\end{equation}
give the relative estimates as follows.  At the regular fan, cyclic
symmetry makes $DF_4(a\mathbf1)$ a constant material covector, so it
vanishes on $\sum_i d_i=0$.  The same symmetry makes
$\Psi_4(a\mathbf1)$ a constant support covector, which vanishes on the
area/translation quotient.  Taylor's formula and \eqref{CT:eq:Cauchy-bounds}
therefore give
\[
 |F_4(\alpha)-F_4(a\mathbf1)|
 \le \omega_4(S)\int_0^1(1-t)E(t)\,dt
 \le C\omega_4(S)\Jfour(\alpha).
\]
Apply the fundamental theorem of calculus to the fixed-fiber covector
$\Psi_4(h(t))\circ\mathcal T_t$.  The regular-fan value vanishes on the
transported quotient, and Cauchy--Schwarz gives
\[
 |\Psi_4(\alpha)[Z]|
 \le \omega_{4,s}(S)\left(\int_0^1E(t)\,dt\right)^{1/2}
       \|Z\|_{\Xjp_\alpha}
 \le C\omega_{4,s}(S)\Jfour(\alpha)^{1/2}\|Z\|_{\Xjp_\alpha}.
\]
These are the pure and mixed tail bounds in the uniformly comparable
regime.  The following elementary dichotomy verifies that this regime and
the absolute regime exhaust all positive fans.  The exact Bregman identity
is
\begin{equation}\label{CT:eq:J4-bregman-comparability}
 \Jfour(\alpha)=\sum_i d_i^2
  (\alpha_i^2+2a\alpha_i+3a^2).
\end{equation}
Let $m=\max_i\alpha_i$, so $S=a+m$, and put
$\rho=\Jfour(\alpha)^{1/2}/S^2$.  Since $m\ge a$ and the summand at an
index where $\alpha_i=m$ is at least $(m-a)^2m^2$,
\[
 \frac{m-a}{S}\le2\rho.
\]
If $\rho\le1/8$, then $a=(S-(m-a))/2\ge3S/8$; applying
\eqref{CT:eq:J4-bregman-comparability} at every index gives
\[
 \max_i\frac{|d_i|}{S}
 \le\frac{8}{3\sqrt3}\rho<2\rho,
 \qquad
 \max_i\frac{|d_i|}{a}\le\frac{16}{3}\rho.
\]
Fix $0<\beta<1$.  If $\Jfour<S^{5-\beta}$, then
$\rho<S^{(1-\beta)/2}$; for all sufficiently small $S$ the last bound is
at most $1/2$, and the whole segment is uniformly comparable by convexity.
In the complementary case $\Jfour\ge S^{5-\beta}$,
\eqref{CT:eq:tail-majorants} gives respectively the relative factors
$CS^\beta$ and $CS^{1+\beta/2}$.  Thus the two cases are exhaustive and
both are $o(1)$ multiples of the required defects.
Thus, under CT--M4, the quartic-and-higher Taylor tail has the uniform
material bounds required by CT.  This deduction uses no arbitrary-order
tensor estimate.

\subsection{Literal fan-to-bubble dictionary}

Choose normal angles $\theta_i$ with
$\theta_{i+1}-\theta_i=\alpha_i$ and set
\begin{equation}\label{CT:eq:cells}
 K_i=[\theta_i-\alpha_{i-1}/2,\theta_i+\alpha_i/2].
\end{equation}
On $K_i$, $x=\theta-\theta_i$, the exact normal fan input is
\begin{equation}\label{CT:eq:secant}
 f_\alpha(\theta)=h_\alpha\sec x-1,
 \qquad
 h_\alpha=\left(\frac{\pi}{\sum_j\tan(\alpha_j/2)}\right)^{1/2}.
\end{equation}
Define, on the same cell and with no change of nodes,
\begin{align}
 c_\alpha&=h_\alpha-1,\notag\\
 q_\alpha(\theta)&=\tfrac12x^2,\label{CT:eq:bubble}\\
 r_\alpha(\theta)&=(h_\alpha-1)\tfrac12x^2
 +h_\alpha(\sec x-1-\tfrac12x^2).\notag
\end{align}
Then the coordinate dictionary is the literal equality
\begin{equation}\label{CT:eq:dictionary}
 \boxed{f_\alpha=c_\alpha+q_\alpha+r_\alpha.}
\end{equation}
The function $q_\alpha$ is continuous at the midpoint vertices because
both adjacent values are $\alpha_i^2/8$, and
\begin{equation}\label{CT:eq:q-distribution}
 q_\alpha''=1-\sum_i\alpha_i
 \delta_{\theta_i+\alpha_i/2}.
\end{equation}
Taylor's formula on each normalized half-cell gives
\begin{equation}\label{CT:eq:r-bounds}
 \|r_\alpha\|_\infty\le CS^4,\qquad
 \|r_\alpha\|_{W^{1,\infty}}\le CS^3,\qquad
 \|r_\alpha\|_{H^{1/2}}^2\le CS^7.
\end{equation}
Constants vanish in \eqref{CT:eq:null}, including after polarization, because
$\int(\Dop g)(\Dop h)=\int g'h'$.  The tame estimate for the null form and
\eqref{CT:eq:r-bounds} place every term containing $r_\alpha$ in the
order-one/tail remainder already controlled above.  The normal component
of the physical support slot is exactly the trace $v_\alpha(z)$ used in
PT.  Therefore the two surviving principal rows of the third derivative
are precisely
\begin{equation}\label{CT:eq:terminal-rows}
 \kappa_{\mathbb D}\Cnull(q_\alpha,q_\alpha,q_\alpha),
 \qquad
 \kappa_{\mathbb D}\Cnull(q_\alpha,q_\alpha,v_\alpha(z)),
\end{equation}
with the same nodes, normalization, and projection as the quadratic-bubble
interface; the common fixed coefficient $\kappa_{\mathbb D}$ remains
outside both terminal estimates.

For downstream Taylor assembly one must use
$\mathcal T_3=(3!)^{-1}D^3\mathcal L(0)$.  Thus the corresponding pure
Taylor row and differentiated mixed Taylor row are
\begin{equation}\label{CT:eq:terminal-taylor-rows}
 \boxed{
 \frac{\kappa_{\mathbb D}}6
   \Cnull(q_\alpha,q_\alpha,q_\alpha),
 \qquad
 \frac{\kappa_{\mathbb D}}2
   \Cnull(q_\alpha,q_\alpha,v_\alpha(z)).}
\end{equation}
The second coefficient is $3(\kappa_{\mathbb D}/6)$, because
differentiating $\mathcal T_3(W,W,W)$ in one support slot produces three
equal placements.

\subsection{CT reduction theorem and audit}

\begin{theorem}[Exact CT core]
\label{CT:thm:CT}
Unconditionally, the fixed-disk third derivative is the sum of
$\kappa_{\mathbb D}$ times the null form \eqref{CT:eq:null} and the explicit
normal order-one ledger \eqref{CT:eq:R-ledger}.  The full
physical-coordinate tensor is \eqref{CT:eq:physical-complete-split}; its area
and translation rows vanish exactly.  The radial null-form outputs are
\eqref{CT:eq:terminal-rows}, with Taylor factors
\eqref{CT:eq:terminal-taylor-rows}; the nonzero curvature support symbol
\eqref{CT:eq:effective-leading-forcing} is retained separately and exported
to the JP--JC audit.  All unconditional stated constants are independent
of the number and minimum size of active cells.  The old material-tail
proposal is not used by the active DAG.
\end{theorem}

\begin{proof}
Combine Proposition~\ref{CT:prop:full-tensor}, Theorem~\ref{CT:thm:complete-split},
Lemma~\ref{CT:lem:R-relative}, Propositions
\ref{CT:prop:curvature-counterexample} and
\ref{CT:prop:joint-mixed-counterexample}, and \eqref{CT:eq:dictionary}.  The
None of these statements imports an estimate for either terminal null-form
row, and the active proof imports no premise from the archived calculation.
\end{proof}

\begin{center}
\begin{longtable}{>{\raggedright\arraybackslash}p{.30\linewidth}
                  >{\raggedright\arraybackslash}p{.58\linewidth}}
\toprule
audit item & disposition\\
\midrule
\endfirsthead
\toprule
audit item & disposition\\
\midrule
\endhead
state rows & all three families displayed in \eqref{CT:eq:full-tensor};\\
metric, mass, and scale & generated exhaustively by
\eqref{CT:eq:J-jets}--\eqref{CT:eq:H-jets};\\
area and translation projection & explicit coefficients
\eqref{CT:eq:chi}, then exact kernel cancellation;\\
curvature/tangential support & both the separate bound and CT--JM are
retracted; \eqref{CT:eq:effective-leading-forcing} is retained for JP--JC;\\
principal symbol & same-sign null form \eqref{CT:eq:null};\\
all lower rows & disjoint ledger \eqref{CT:eq:R-ledger} and
order-one estimate \eqref{CT:eq:R-estimate};\\
quartic pure/linear tail & old CT--M4 proposal and its consequences are
retained only as an archive; the active exact-support route uses neither;\\
bubble output & literal equality \eqref{CT:eq:dictionary}, with identical
cells and midpoint atoms.\\
\bottomrule
\end{longtable}
\end{center}
\endgroup
